% Draft replacement for lines 1357–1369
% Penalized form, empirical, footnote on Lagrangian equivalence

The results above establish that balancing weights can beat the true weights, but they do not tell us by how much. These results are part of a long tradition of work on minimax estimation of linear functionals \citep{ibragimov1985nonparametric, donoho1991geometrizing, donoho1994statistical, armstrong2018optimal} that provides the tools to answer this. That literature studies problems like ours: estimating a linear functional $\psi_1 = \E[\mu_1(X)]$ when $\mu_1$ is known to lie in a convex function class. Its central object is the \emph{modulus of continuity},\footnote{This is the Lagrangian dual of the constrained modulus $\bar\omega(\delta) = \sup\{|\frac{1}{n}\sum_i m_1(X_i) - \frac{1}{n}\sum_i m_2(X_i)| : m_1, m_2 \in \model,\, \norm{m_1 - m_2}_{L_2(\Pn)} \le \delta\}$ used by \citet{donoho1994statistical} and \citet{armstrong2018optimal}.}
\begin{equation}
\label{eq:modulus-of-continuity}
\omega(s) = \sup_{m \in \model} \left\{ \frac{1}{n}\sum_{i=1}^n m(X_i) - \frac{1}{2sn} \sum_{W_i=1} m(X_i)^2 \right\},
\end{equation}
which measures how much $\frac{1}{n}\sum_i m(X_i)$ can vary over $m \in \model$, net of a quadratic penalty on $m$'s empirical norm on the treated sample. The maximal imbalance of the minimax balancing weights has a clean characterization in terms of $\omega$. At the optimal bias-variance tradeoff parameterized by $s > 0$,
\begin{equation}
\label{eq:donoho-liu-bias}
\imbalance_\model(\hgamma) = \omega(s) - \omega'(s)\,s.
\end{equation}
When $\omega$ is approximately linear---$\omega(s) \approx c\,s$ for some constant $c$---the maximal imbalance is approximately zero. So bounding imbalance reduces to bounding the deviation of $\omega$ from linearity.
